\chapter{Conception du Système}
\label{chap:conception}

\section{Introduction}

Ce chapitre présente la phase de conception de l'application développée pour CIVCO BTP.
Il traduit les besoins métier identifiés au Chapitre~1 en spécifications techniques
formelles et en modèles architecturaux. Nous commençons par identifier les acteurs et
leurs cas d'utilisation, puis nous définissons les exigences fonctionnelles et
non-fonctionnelles du système. Nous décrivons ensuite l'architecture retenue, la
conception des modules métier principaux, et nous terminons par la modélisation
\textbf{UML} : diagrammes de cas d'utilisation, de séquence, de classes et d'activité.

% ─────────────────────────────────────────────────────────────────────
\section{Analyse des besoins}

\subsection{Identification des acteurs}

Six catégories d'acteurs interagissent avec l'application de gestion de CIVCO BTP :

\begin{description}
  \item[Administrateur] Responsable de la configuration globale de l'application :
  gestion des utilisateurs, des rôles et permissions, des paramètres documentaires,
  du thème visuel et de la consultation de l'historique des actions.
  \item[Chef de projet] Pilote un ou plusieurs chantiers : suivi des phases et des
  tâches, affectation des équipes, consultation des documents liés au projet,
  gestion des contrats et des avenants.
  \item[Commercial] Gère la relation client et les documents commerciaux : création
  et suivi des devis, conversion en factures, production des bons de livraison.
  \item[Comptable] Assure le suivi financier : consultation et validation des
  factures, suivi des paiements, état de facturation par projet.
  \item[Utilisateur interne] Collaborateur terrain ou technique assigné à des tâches
  spécifiques : consultation de ses missions, mise à jour de l'avancement, accès
  limité aux projets auxquels il est affecté.
  \item[Client externe] Maître d'ouvrage ou client final accédant au portail dédié
  pour consulter l'état de ses projets, accepter les devis, signer les contrats et
  échanger avec l'équipe CIVCO BTP.
\end{description}

\subsection{Cas d'utilisation}

Les cas d'utilisation principaux, regroupés par domaine fonctionnel, sont les suivants :

\textbf{Authentification et tableau de bord :}
\begin{itemize}
  \item UC-01 : S'authentifier à l'application (connexion sécurisée par session).
  \item UC-02 : Consulter le tableau de bord synthétique (projets en cours, tâches,
  indicateurs clés).
\end{itemize}

\textbf{Gestion des projets :}
\begin{itemize}
  \item UC-03 : Créer et modifier un projet de construction (référence, client, budget,
  dates, localisation).
  \item UC-04 : Définir les phases d'un projet et suivre l'avancement global.
  \item UC-05 : Affecter des membres d'équipe à un projet.
  \item UC-06 : Consulter la carte géographique des chantiers.
\end{itemize}

\textbf{Gestion des clients :}
\begin{itemize}
  \item UC-07 : Créer, modifier et archiver une fiche client.
  \item UC-08 : Gérer les contacts associés à un client.
  \item UC-09 : Provisionner un accès portail pour un client externe.
\end{itemize}

\textbf{Documents commerciaux :}
\begin{itemize}
  \item UC-10 : Créer un devis avec lignes de détail et le lier à un projet.
  \item UC-11 : Convertir un devis accepté en facture.
  \item UC-12 : Émettre un bon de livraison (\textbf{BL}) lié à un projet.
  \item UC-13 : Imprimer ou exporter un document commercial (devis, facture, BL).
\end{itemize}

\textbf{Contrats et tâches :}
\begin{itemize}
  \item UC-14 : Créer et suivre un contrat associé à un projet.
  \item UC-15 : Enregistrer un avenant contractuel (modification budget/durée).
  \item UC-16 : Créer, assigner et suivre des tâches opérationnelles.
\end{itemize}

\textbf{Administration et portail client :}
\begin{itemize}
  \item UC-17 : Configurer les rôles, permissions et paramètres de l'application.
  \item UC-18 : Consulter le journal d'activité et l'historique des actions.
  \item UC-19 : Échanger des messages avec un client via la messagerie intégrée.
  \item UC-20 : (Client externe) Consulter ses projets, accepter un devis, signer
  un contrat via le portail.
\end{itemize}

\subsection{Exigences fonctionnelles}

\begin{itemize}
  \item \textbf{EF-01 Authentification :} Connexion par email et mot de passe via
  Laravel Sanctum (session sécurisée par cookie, protection \textbf{CSRF}).
  \item \textbf{EF-02 Gestion des rôles :} Attribution de rôles métier (administrateur,
  chef de projet, commercial, comptable, utilisateur) avec permissions granulaires
  (\textbf{RBAC}).
  \item \textbf{EF-03 Gestion des projets :} \textbf{CRUD} complet sur les projets,
  avec phases, équipe, budget, avancement, géolocalisation du chantier.
  \item \textbf{EF-04 Gestion des clients :} Fiches clients, contacts, archivage,
  provisionnement d'accès portail.
  \item \textbf{EF-05 Devis :} Création de devis multi-lignes, calcul automatique des
  totaux HT/TVA/TTC, suivi des statuts (brouillon, envoyé, accepté, refusé).
  \item \textbf{EF-06 Factures :} Conversion devis $\rightarrow$ facture, suivi de
  l'état de paiement, numérotation automatique.
  \item \textbf{EF-07 Bons de livraison :} Production de BL liés aux projets, avec
  lignes de détail et suivi des quantités livrées.
  \item \textbf{EF-08 Contrats :} Gestion des contrats et avenants, compilation à
  partir de modèles, double signature (entreprise + client).
  \item \textbf{EF-09 Tâches :} Planification par phase, assignation, statuts
  (à faire, en cours, terminé), filtrage par utilisateur.
  \item \textbf{EF-10 Tableau de bord :} Widgets configurables (projets actifs,
  tâches prioritaires, graphiques d'avancement).
  \item \textbf{EF-11 Portail client :} Interface dédiée pour consultation des projets,
  acceptation de devis et signature de contrats.
  \item \textbf{EF-12 Messagerie :} Échanges projet-scopés entre équipe interne et
  client externe.
  \item \textbf{EF-13 Historique :} Journalisation des actions sensibles (création,
  modification, archivage) consultable par l'administrateur.
\end{itemize}

\subsection{Exigences non-fonctionnelles}

\textbf{Performance :}
\begin{itemize}
  \item \textbf{ENF-01 :} Temps de réponse des API inférieur à 500~ms pour les
  opérations de lecture courantes (listes, fiches détail).
  \item \textbf{ENF-02 :} Chargement initial de l'interface (SPA) inférieur à 3~s
  sur une connexion standard.
\end{itemize}

\textbf{Sécurité :}
\begin{itemize}
  \item \textbf{ENF-03 :} Mots de passe hachés (bcrypt) ; sessions Sanctum avec
  expiration configurable.
  \item \textbf{ENF-04 :} Contrôle d'accès par permission sur chaque route API ;
  vérification côté serveur avant toute action métier.
  \item \textbf{ENF-05 :} Protection \textbf{CSRF} sur les requêtes mutantes ;
  communication \textbf{HTTPS} en production.
\end{itemize}

\textbf{Ergonomie et accessibilité :}
\begin{itemize}
  \item \textbf{ENF-06 :} Interface responsive adaptée aux postes de bureau et
  tablettes de chantier.
  \item \textbf{ENF-07 :} Support bilingue français / anglais via fichiers de
  traduction centralisés.
\end{itemize}

\textbf{Maintenabilité :}
\begin{itemize}
  \item \textbf{ENF-08 :} Architecture modulaire : contrôleurs fins, services métier,
  modèles Eloquent, composants React par feature.
  \item \textbf{ENF-09 :} Schéma de base versionné par migrations Laravel,
  reproductible entre environnements.
\end{itemize}

\subsection{Scénarios d'utilisation détaillés}

Pour illustrer concrètement les exigences fonctionnelles, trois scénarios types ont
été formalisés lors de la phase de conception :

\medskip

\textbf{Scénario 1 --- Lancement d'un nouveau chantier :}
un chef de projet crée un projet en sélectionnant le client, renseigne le budget et
les dates, définit les phases (études, terrassement, gros œuvre, finitions), affecte
l'équipe (ingénieur, conducteur de travaux, technicien) et positionne le chantier
sur la carte. Le commercial rattache un devis préalablement validé.

\medskip

\textbf{Scénario 2 --- Cycle commercial complet :}
le commercial crée un devis multi-lignes pour un projet existant, l'envoie au client
via le portail, le client l'accepte en ligne, le comptable convertit le devis en
facture, enregistre un paiement partiel puis le solde. Chaque étape est tracée dans
le journal d'activité.

\medskip

\textbf{Scénario 3 --- Suivi opérationnel quotidien :}
un conducteur de travaux consulte ses tâches du jour sur tablette, met à jour le
statut d'une tâche terminée, le chef de projet ajuste l'avancement global du projet
et le tableau de bord reflète immédiatement la progression.

% ─────────────────────────────────────────────────────────────────────
\section{Architecture du système}

\subsection{Architecture globale en trois couches}

L'application adopte une architecture trois couches, avec une séparation nette entre
présentation, logique métier et persistance :

\begin{tcolorbox}[colback=ensamBlue!5, colframe=ensamBlue,
  title={\textbf{Couche Présentation} : React 19 + Vite + Tailwind CSS}, arc=3pt, boxrule=0.8pt]
Interface web organisée en \textbf{SPA} : tableau de bord, modules projets, clients,
devis, factures, tâches, configuration et portail client. Navigation par React Router,
état global via Context API (authentification, thème, langue).
\end{tcolorbox}

\vspace{-0.2cm}

\begin{tcolorbox}[colback=ensamBlue!10, colframe=ensamBlue,
  title={\textbf{Couche Logique métier} : Laravel 13 + Sanctum}, arc=3pt, boxrule=0.8pt]
API \textbf{REST} sous \texttt{/api/v1/} : contrôleurs par domaine, services métier,
validation (Form Requests), transformation JSON (API Resources), middleware de
permissions et journalisation d'activité.
\end{tcolorbox}

\vspace{-0.2cm}

\begin{tcolorbox}[colback=ensamBlue!15, colframe=ensamBlue,
  title={\textbf{Couche Persistance} : MySQL / PostgreSQL + Eloquent ORM}, arc=3pt, boxrule=0.8pt]
Stockage relationnel de toutes les entités métier : utilisateurs, rôles, clients,
projets, documents commerciaux, contrats, tâches. Migrations versionnées et seeders
de démonstration.
\end{tcolorbox}

Les échanges entre ces couches s'effectuent en \textbf{JSON} via \textbf{HTTPS} côté
présentation, et en \textbf{SQL} via l'\textbf{ORM} Eloquent côté persistance.

\subsection{Organisation du backend Laravel}

Le backend est structuré selon le patron \textbf{MVC} étendu par une couche de services :

\begin{itemize}
  \item \textbf{Contrôleurs} (\texttt{app/Http/Controllers/Api/V1/}) : points d'entrée
  HTTP, délégation aux services, retour de Resources JSON.
  \item \textbf{Form Requests} (\texttt{app/Http/Requests/}) : validation des données
  entrantes par domaine (StoreProjectRequest, StoreQuoteRequest, etc.).
  \item \textbf{Services} (\texttt{app/Services/}) : logique métier complexe
  (QuoteService, InvoiceService, DocumentService, DashboardService, etc.).
  \item \textbf{Modèles} (\texttt{app/Models/}) : entités Eloquent avec relations,
  casts et scopes.
  \item \textbf{Resources} (\texttt{app/Http/Resources/}) : sérialisation JSON
  normalisée pour le frontend.
  \item \textbf{Middleware} : authentification Sanctum, vérification des permissions,
  résolution du contexte entreprise.
\end{itemize}

\subsection{Organisation du frontend React}

Le frontend suit une organisation par \textbf{features} (domaines métier) :

\begin{itemize}
  \item \texttt{features/projects/} : liste, fiche détail, onglets phases/tâches/équipe.
  \item \texttt{features/clients/} : gestion clients, portail, archivage.
  \item \texttt{features/quotes/}, \texttt{invoices/}, \texttt{delivery-forms/} :
  documents commerciaux.
  \item \texttt{features/tasks/} : planification et suivi des tâches.
  \item \texttt{features/dashboard/} : tableau de bord et widgets.
  \item \texttt{features/configuration/} : paramètres, rôles, modèles de documents.
  \item \texttt{api/} : modules client HTTP (Axios) par ressource API.
\end{itemize}

\subsection{Protocoles de communication}

L'application repose principalement sur le protocole \textbf{REST} (HTTP/JSON) :

\begin{itemize}
  \item \textbf{Authentification} : flux Sanctum SPA (cookie de session + jeton CSRF).
  \item \textbf{Opérations CRUD} : verbes HTTP standard sur les ressources métier
  (\texttt{GET /api/v1/projects}, \texttt{POST /api/v1/quotes}, etc.).
  \item \textbf{Upload de fichiers} : requêtes \texttt{multipart/form-data} pour les
  documents joints, médias de projet et signatures.
  \item \textbf{Impression} : génération de PDF côté client (composants React dédiés)
  avec suivi des impressions officielles.
\end{itemize}

\subsection{Flux d'authentification}

Le flux d'authentification Sanctum pour une SPA suit les étapes suivantes :

\begin{enumerate}
  \item Le frontend appelle \texttt{GET /sanctum/csrf-cookie} pour obtenir un jeton
  \textbf{CSRF}.
  \item L'utilisateur saisit ses identifiants ; le frontend envoie
  \texttt{POST /api/v1/login} avec le jeton CSRF.
  \item Laravel valide les identifiants, crée une session et retourne le profil
  utilisateur (rôles, permissions).
  \item Les requêtes suivantes incluent automatiquement le cookie de session ;
  chaque route protégée vérifie l'authentification et les permissions.
  \item En cas de session expirée (\texttt{401}), le frontend redirige vers la page
  de connexion.
\end{enumerate}

\subsection{Matrice des permissions par profil}

Le Tableau~\ref{tab:permissions} présente une synthèse des droits d'accès par profil
métier. Les permissions exactes sont configurables par l'administrateur ; le tableau
reflète la configuration type déployée pour CIVCO BTP.

\begin{table}[H]
\caption{Matrice des permissions par profil utilisateur (synthèse)}
\label{tab:permissions}
\centering
\small
\begin{tabular}{|l|c|c|c|c|c|}
\hline
\textbf{Module / Action} & \textbf{Admin} & \textbf{Chef proj.} & \textbf{Com.} &
\textbf{Compt.} & \textbf{Utilis.} \\
\hline
Tableau de bord & $\checkmark$ & $\checkmark$ & $\checkmark$ & $\checkmark$ & $\checkmark$ \\
\hline
Projets (lecture) & $\checkmark$ & $\checkmark$ & $\checkmark$ & $\checkmark$ & $\checkmark$* \\
\hline
Projets (écriture) & $\checkmark$ & $\checkmark$ & --- & --- & --- \\
\hline
Clients & $\checkmark$ & $\checkmark$ & $\checkmark$ & --- & --- \\
\hline
Devis & $\checkmark$ & $\checkmark$ & $\checkmark$ & --- & --- \\
\hline
Factures & $\checkmark$ & $\checkmark$ & $\checkmark$ & $\checkmark$ & --- \\
\hline
Tâches (toutes) & $\checkmark$ & $\checkmark$ & --- & --- & --- \\
\hline
Tâches (propres) & $\checkmark$ & $\checkmark$ & $\checkmark$ & $\checkmark$ & $\checkmark$ \\
\hline
Configuration & $\checkmark$ & --- & --- & --- & --- \\
\hline
Historique & $\checkmark$ & --- & --- & --- & --- \\
\hline
\end{tabular}

\vspace{0.2cm}
{\footnotesize * Accès limité aux projets assignés.}
\end{table}

% ─────────────────────────────────────────────────────────────────────
\section{Conception des modules métier}

Cette section détaille la conception fonctionnelle des principaux modules de
l'application, en précisant les entités manipulées, les règles de gestion et les
interactions entre modules.

\subsection{Module Projets}

Le module Projets constitue le cœur de l'application. Un projet représente un chantier
ou une affaire de construction, lié à un client, découpé en phases et suivi par une
équipe interne.

\medskip

\textbf{Entités principales :} \texttt{Project}, \texttt{ProjectPhase}, \texttt{Expense},
\texttt{ProjectMedia}, \texttt{ProgressSnapshot}.

\medskip

\textbf{Règles de gestion :}
\begin{itemize}
  \item Chaque projet possède une référence unique et un statut (planifié, en cours,
  suspendu, terminé, archivé).
  \item L'avancement (\texttt{progress\_percent}) peut être mis à jour manuellement
  ou calculé à partir des tâches terminées.
  \item La géolocalisation (latitude/longitude) permet l'affichage sur la carte des
  chantiers (Leaflet).
  \item Seuls les utilisateurs disposant de la permission \texttt{project.manage}
  peuvent créer ou modifier un projet.
\end{itemize}

\subsection{Module Clients}

Le module Clients centralise les informations des maîtres d'ouvrage et clients
commerciaux de CIVCO BTP.

\medskip

\textbf{Entités principales :} \texttt{Client}, \texttt{ClientContact}, \texttt{Badge}.

\medskip

\textbf{Règles de gestion :}
\begin{itemize}
  \item Un client peut être archivé sans suppression des projets associés (soft archive).
  \item La création d'un accès portail génère un compte utilisateur lié au client
  (\texttt{client\_id} sur \texttt{User}).
  \item Les contacts multiples sont gérés via la relation \texttt{ClientContact}.
\end{itemize}

\subsection{Module Documents commerciaux}

Ce module regroupe les devis, factures et bons de livraison, qui partagent une
structure commune de lignes de détail.

\medskip

\textbf{Entités principales :} \texttt{Quote}, \texttt{QuoteLine}, \texttt{Invoice},
\texttt{InvoiceLine}, \texttt{DeliveryForm}, \texttt{DeliveryFormLine}, \texttt{Payment}.

\medskip

\textbf{Flux documentaire :}
\begin{enumerate}
  \item Le commercial crée un devis (statut \textit{brouillon}) lié à un projet.
  \item Le devis est validé et envoyé au client (statut \textit{envoyé}).
  \item Le client accepte via le portail ou l'équipe interne confirme l'acceptation.
  \item Le devis est converti en facture (copie des lignes, nouvelle référence).
  \item Les paiements partiels ou totaux sont enregistrés sur la facture.
  \item Les bons de livraison documentent les livraisons matérielles sur chantier.
\end{enumerate}

\subsection{Module Contrats}

Le module Contrats permet de formaliser l'engagement contractuel lié à un projet,
avec support des avenants.

\medskip

\textbf{Entités principales :} \texttt{Contract}, \texttt{ContractAmendment},
\texttt{ContractTemplate}, \texttt{DocumentTemplate}.

\medskip

\textbf{Règles de gestion :}
\begin{itemize}
  \item Un contrat est généré à partir d'un modèle (\texttt{ContractTemplate}) et
  des données du projet (compilation de champs dynamiques).
  \item Les avenants modifient le budget ou la durée sans recréer le contrat initial.
  \item La signature client est capturée via un pad de signature sur le portail.
\end{itemize}

\subsection{Module Tâches}

Le module Tâches organise le travail opérationnel au sein de chaque projet.

\medskip

\textbf{Entités principales :} \texttt{Task} (liée à \texttt{Project} et optionnellement
à \texttt{ProjectPhase}).

\medskip

\textbf{Règles de gestion :}
\begin{itemize}
  \item Une tâche est assignée à un utilisateur et possède un statut, une priorité
  et une date d'échéance.
  \item Les permissions \texttt{task.view\_own} et \texttt{task.view\_all} contrôlent
  la visibilité selon le profil.
  \item Le tableau de bord agrège les tâches en retard et à venir.
\end{itemize}

\subsection{Module Tableau de bord et carte}

Le tableau de bord offre une vue consolidée de l'activité de CIVCO BTP :

\begin{itemize}
  \item Nombre de projets actifs, en retard, terminés.
  \item Graphiques d'avancement (Recharts).
  \item Widget des tâches prioritaires de l'utilisateur connecté.
  \item Carte Leaflet des projets géolocalisés.
\end{itemize}

Les widgets sont conditionnés par les permissions de l'utilisateur : un commercial
voit les indicateurs commerciaux, un chef de projet les indicateurs de chantier.

\subsection{Module Portail client et messagerie}

Le portail client constitue l'interface dédiée aux clients externes de CIVCO BTP.
Accessible via une route séparée (\texttt{/portal}), il offre une expérience simplifiée
centrée sur la consultation et la validation.

\medskip

\textbf{Fonctionnalités portail :}
\begin{itemize}
  \item Tableau de bord client : projets en cours, documents en attente de signature.
  \item Consultation des devis et acceptation en ligne.
  \item Signature électronique des contrats (pad de signature canvas).
  \item Fil de discussion avec l'équipe projet (si autorisé).
  \item Calendrier des échéances et jalons du projet.
\end{itemize}

\medskip

\textbf{Messagerie interne :} Les échanges entre l'équipe CIVCO BTP et le client sont
scopés par projet. Seuls les membres d'équipe disposant de l'autorisation
\texttt{can\_chat\_with\_client} peuvent initier ou répondre à une conversation.
Les messages sont persistés en base (\texttt{PortalMessage}) et consultables depuis
l'interface d'administration.

\subsection{Module Configuration et historique}

Le module Configuration regroupe les paramètres transverses de l'application :

\begin{itemize}
  \item \textbf{Rôles et permissions} : création de profils métier, attribution de
  permissions granulaires (\texttt{project.view}, \texttt{quote.manage}, etc.).
  \item \textbf{Modèles de documents} : templates de contrats et documents commerciaux
  avec champs dynamiques (placeholders).
  \item \textbf{Contrôles d'impression} : quotas d'impressions officielles, filigrane
  pour les copies non officielles.
  \item \textbf{Thème visuel} : couleurs de l'interface, logo de l'entreprise.
  \item \textbf{Référentiels} : badges clients, lots, secteurs d'activité.
\end{itemize}

\medskip

Le module Historique (\texttt{ActivityLog}) enregistre les actions sensibles :
création/modification de projets, validation de devis, archivage de clients.
Chaque entrée contient l'utilisateur, l'action, l'entité concernée et l'horodatage,
offrant une traçabilité utile à l'administration et à l'audit interne.

% ─────────────────────────────────────────────────────────────────────
\section{Modélisation de la base de données}

\subsection{Principes de conception}

La base de données relationnelle constitue le socle de persistance de l'application.
Sa conception repose sur les principes suivants :

\begin{itemize}
  \item \textbf{Normalisation} : les entités sont décomposées pour éviter la
  redondance (clients séparés des projets, lignes de devis dans une table dédiée).
  \item \textbf{Intégrité référentielle} : les clés étrangères garantissent la
  cohérence (un devis appartient toujours à un projet et un client existants).
  \item \textbf{Versionnement} : le schéma évolue via des migrations Laravel,
  reproductibles entre environnements de développement, test et production.
  \item \textbf{Traçabilité} : les horodatages (\texttt{created\_at},
  \texttt{updated\_at}) et le journal d'activité complètent l'historique métier.
\end{itemize}

\subsection{Tables principales}

Le Tableau~\ref{tab:db_tables} présente les tables centrales du modèle de données
et leur rôle dans l'application.

\begin{table}[H]
\caption{Tables principales de la base de données}
\label{tab:db_tables}
\centering
\small
\begin{tabular}{|l|p{4.5cm}|p{5.5cm}|}
\hline
\textbf{Table} & \textbf{Entité} & \textbf{Rôle} \\
\hline
\texttt{users} & Utilisateur & Comptes internes et clients portail \\
\hline
\texttt{roles}, \texttt{permissions} & Sécurité & Profils métier et droits granulaires \\
\hline
\texttt{clients} & Client & Maîtres d'ouvrage et clients commerciaux \\
\hline
\texttt{projects} & Projet & Chantiers, budget, avancement, coordonnées \\
\hline
\texttt{project\_phases} & Phase & Découpage temporel d'un projet \\
\hline
\texttt{tasks} & Tâche & Missions opérationnelles assignées \\
\hline
\texttt{quotes}, \texttt{quote\_lines} & Devis & Documents commerciaux et lignes \\
\hline
\texttt{invoices}, \texttt{invoice\_lines} & Facture & Facturation et détail des prestations \\
\hline
\texttt{contracts} & Contrat & Engagements contractuels liés aux projets \\
\hline
\texttt{activity\_logs} & Journal & Traçabilité des actions sensibles \\
\hline
\end{tabular}
\end{table}

Au total, le schéma compte plus de soixante-dix tables lorsque l'on inclut les
tables de liaison (équipe projet, badges clients, paiements, impressions, etc.),
les référentiels (lots, secteurs) et les tables techniques (sessions, notifications).

\subsection{Relations et contraintes d'intégrité}

Les relations cardinales les plus importantes structurent le modèle métier :

\begin{itemize}
  \item Un \textbf{client} possède zéro à plusieurs \textbf{projets} ; un projet
  est rattaché à un seul client.
  \item Un \textbf{projet} regroupe plusieurs \textbf{phases}, \textbf{tâches},
  \textbf{devis}, \textbf{factures} et \textbf{contrats}.
  \item Un \textbf{devis accepté} peut être converti en une seule \textbf{facture}
  (relation 1:0..1), les lignes étant recopiées automatiquement.
  \item Un \textbf{utilisateur} peut appartenir à l'équipe de plusieurs projets
  (relation N:N via \texttt{project\_user}).
  \item Un \textbf{rôle} agrège plusieurs \textbf{permissions} ; chaque utilisateur
  interne est associé à un rôle métier.
\end{itemize}

Les contraintes \texttt{ON DELETE RESTRICT} sur les clés étrangères critiques
(client, projet) empêchent la suppression accidentelle d'enregistrements encore
référencés. L'archivage des clients et la clôture des projets sont privilégiés
à la suppression physique.

% ─────────────────────────────────────────────────────────────────────
\section{Conception de l'interface utilisateur}

\subsection{Principes ergonomiques}

L'interface a été conçue pour un usage professionnel quotidien, en tenant compte
des contraintes des utilisateurs de CIVCO BTP :

\begin{itemize}
  \item \textbf{Lisibilité} : thème sombre par défaut, contrastes élevés, typographie
  claire pour une utilisation prolongée devant écran.
  \item \textbf{Efficacité} : accès rapide aux modules via une barre latérale fixe,
  recherche globale et raccourcis depuis le tableau de bord.
  \item \textbf{Cohérence} : chaque module suit le même patron (liste $\rightarrow$
  fiche détail $\rightarrow$ actions contextuelles).
  \item \textbf{Responsive} : adaptation aux tablettes utilisées sur chantier pour
  la consultation des tâches et des documents.
  \item \textbf{Feedback utilisateur} : messages de confirmation, indicateurs de
  chargement et alertes en cas d'erreur de validation.
\end{itemize}

\subsection{Charte graphique et navigation}

L'identité visuelle de CIVCO BTP est intégrée via un système de thème dynamique :
couleur principale (\texttt{ensamBlue} / bleu entreprise), logo en en-tête et
variables CSS injectées depuis la configuration. La navigation repose sur une
\textbf{sidebar} persistante dont les entrées sont filtrées selon les permissions
de l'utilisateur connecté.

\medskip

La structure de navigation distingue trois zones :

\begin{enumerate}
  \item \textbf{Zone opérationnelle} : tableau de bord, projets, clients, tâches,
  carte des chantiers.
  \item \textbf{Zone documentaire} : devis, factures, bons de livraison.
  \item \textbf{Zone administration} : configuration, historique, gestion d'équipe
  (réservée aux profils autorisés).
\end{enumerate}

Le portail client adopte une interface simplifiée, sans la sidebar interne, centrée
sur la consultation et la validation des documents.

\subsection{Parcours utilisateurs types}

Trois parcours types ont guidé la conception des écrans :

\medskip

\textbf{Parcours commercial :} connexion $\rightarrow$ création d'un client
$\rightarrow$ création d'un projet $\rightarrow$ production d'un devis
$\rightarrow$ envoi au client $\rightarrow$ conversion en facture après acceptation.

\medskip

\textbf{Parcours chef de projet :} connexion $\rightarrow$ consultation du
tableau de bord $\rightarrow$ ouverture d'un chantier $\rightarrow$ mise à jour
des phases et des tâches $\rightarrow$ affectation des collaborateurs
$\rightarrow$ suivi de l'avancement sur la carte.

\medskip

\textbf{Parcours client externe :} connexion portail $\rightarrow$ consultation
des projets en cours $\rightarrow$ acceptation d'un devis $\rightarrow$ signature
électronique du contrat $\rightarrow$ échange via la messagerie intégrée.

% ─────────────────────────────────────────────────────────────────────
\section{Conception de la sécurité et de la traçabilité}

\subsection{Modèle de menaces et contre-mesures}

La conception de la sécurité s'appuie sur une analyse des risques les plus probables
dans le contexte d'une application métier interne :

\begin{table}[H]
\caption{Menaces identifiées et contre-mesures associées}
\label{tab:security}
\centering
\small
\begin{tabular}{|p{4cm}|p{4.5cm}|p{5cm}|}
\hline
\textbf{Menace} & \textbf{Impact} & \textbf{Contre-mesure} \\
\hline
Accès non autorisé & Consultation ou modification de données sensibles & Sanctum, RBAC, middleware permissions \\
\hline
Falsification de requête (CSRF) & Action non voulue au nom d'un utilisateur & Jeton CSRF sur requêtes mutantes \\
\hline
Élévation de privilèges & Contournement des restrictions métier & Vérification côté serveur sur chaque route \\
\hline
Fuite de données client & Atteinte à la confidentialité & Filtrage par projet et par rôle \\
\hline
Perte de traçabilité & Impossibilité d'auditer les actions & Journal d'activité (\texttt{ActivityLog}) \\
\hline
\end{tabular}
\end{table}

\subsection{Journalisation des actions}

Le module \texttt{ActivityLogService} enregistre les opérations sensibles :
création ou modification d'un projet, validation d'un devis, archivage d'un client,
changement de rôle. Chaque entrée comprend l'identifiant de l'utilisateur, le type
d'action, l'entité concernée, l'horodatage et éventuellement les valeurs modifiées.
L'administrateur consulte ce journal depuis le module Historique.

\subsection{Gestion des impressions officielles}

Les documents commerciaux (devis, factures, bons de livraison) peuvent être
imprimés en version \textbf{officielle} (avec en-tête CIVCO BTP) ou en
\textbf{copie non officielle}. Un quota configurable limite le nombre d'impressions
officielles par type de document. Au-delà du quota, le système applique
automatiquement un filigrane « COPIE - NON OFFICIEL » et enregistre chaque
impression dans la table de suivi (\texttt{print\_trackings}).

% ─────────────────────────────────────────────────────────────────────
\section{Modélisation UML}

\subsection{Diagramme de cas d'utilisation}

Le diagramme de cas d'utilisation (Figure~\ref{fig:uc_diagram}) représente les interactions entre les six acteurs
identifiés et le système. L'administrateur dispose des droits les plus étendus
(configuration, historique, gestion des utilisateurs). Le client externe accède
uniquement au portail pour consulter ses projets et valider les documents qui le
concernent.

\medskip

Les relations entre acteurs et cas d'utilisation suivent la logique métier de
CIVCO BTP : le commercial intervient sur les clients et les devis, le comptable
sur les factures, le chef de projet sur les chantiers et les tâches. Les cas
d'utilisation UC-01 à UC-20 détaillés en section~1.2 constituent la base de ce
diagramme. Les cas d'utilisation \textit{include} (par exemple, « S'authentifier »
précède toute action protégée) et \textit{extend} (par exemple, « Imprimer un
document » en complément de « Consulter un devis ») structurent les dépendances
entre fonctionnalités.

\begin{figure}[H]
  \centering
  \includegraphics[width=\textwidth]{usecaseUML.png}
  \caption{Diagramme de cas d'utilisation de l'application CIVCO BTP}
  \label{fig:uc_diagram}
\end{figure}

\subsection{Diagramme de séquence}

Le diagramme de séquence (Figure~\ref{fig:seq_diagram}) illustre le scénario de création d'un devis par un
commercial. Les interactions mettent en évidence la séparation entre l'interface
React, l'API Laravel (validation, service métier) et la persistance en base de
données. Les étapes principales sont les suivantes :
\begin{enumerate}
  \item Le commercial s'authentifie via la SPA (cookie Sanctum + jeton CSRF).
  \item Il sélectionne un projet et saisit les lignes du devis.
  \item Le frontend envoie \texttt{POST /api/v1/quotes} avec les données JSON.
  \item Laravel valide la requête (Form Request), persiste le devis et ses lignes.
  \item L'API retourne une \texttt{QuoteResource} ; l'interface affiche la confirmation.
\end{enumerate}
Ce flux est représentatif du patron adopté sur l'ensemble des modules : le frontend
ne contient aucune logique métier critique ; la validation, les calculs et la
persistance sont toujours assurés côté serveur.

\medskip

D'autres scénarios de séquence ont été modélisés selon le même principe :
\textbf{conversion devis $\rightarrow$ facture} (contrôle du statut, transaction
base de données, copie des lignes), \textbf{signature de contrat} (génération
du PDF, capture de la signature canvas, mise à jour du statut) et
\textbf{assignation de tâche} (notification de l'utilisateur concerné).

\begin{figure}[H]
  \centering
  \includegraphics[width=\textwidth]{sequenceUML.png}
  \caption{Diagramme de séquence : création et enregistrement d'un devis}
  \label{fig:seq_diagram}
\end{figure}

\subsection{Diagramme de classes}

Le diagramme de classes (Figure~\ref{fig:class_diagram}) représente le modèle de données persistant simplifié.
Le Tableau~\ref{tab:classes} complète la description des entités principales et
de leurs relations.

\begin{figure}[H]
  \centering
  \includegraphics[width=\textwidth]{classeUML.png}
  \caption{Diagramme de classes simplifié du modèle de données métier}
  \label{fig:class_diagram}
\end{figure}

\begin{table}[H]
\caption{Description des classes et de leurs relations}
\label{tab:classes}
\centering
\small
\begin{tabular}{|l|p{4.8cm}|p{5.2cm}|}
\hline
\textbf{Classe} & \textbf{Attributs principaux} & \textbf{Relations} \\
\hline
\texttt{User} & id, name, email, role\_id & N User $\rightarrow$ 1 Role ;
N $\leftrightarrow$ N Project (équipe) \\
\hline
\texttt{Role} & id, name, slug & 1 Role $\rightarrow$ N Permission \\
\hline
\texttt{Client} & id, name, email, status & 1 Client $\rightarrow$ N Project ;
1 $\rightarrow$ N ClientContact \\
\hline
\texttt{Project} & id, client\_id, title, status, budget, progress & 1 $\rightarrow$ N Phase,
Task, Quote, Contract \\
\hline
\texttt{Quote} & id, project\_id, reference, total\_ttc, status & 1 $\rightarrow$ N QuoteLine ;
1 $\rightarrow$ 0..1 Invoice \\
\hline
\texttt{Invoice} & id, quote\_id, reference, total\_ttc & 1 $\rightarrow$ N InvoiceLine,
Payment \\
\hline
\texttt{Task} & id, project\_id, title, status, due\_date & N $\rightarrow$ 1 Project ;
N $\rightarrow$ 1 User (assigné) \\
\hline
\texttt{Contract} & id, project\_id, status & 1 $\rightarrow$ N ContractAmendment \\
\hline
\end{tabular}
\end{table}

\medskip

Le modèle met en évidence la \textbf{centralité du projet} : la plupart des entités
métier gravitent autour de \texttt{Project}, ce qui reflète l'organisation interne
de CIVCO BTP où le chantier constitue l'unité de travail de référence. Les
énumérations (\texttt{ProjectStatus}, \texttt{QuoteStatus}, \texttt{InvoiceStatus})
garantissent la cohérence des états dans l'application.

\subsection{Diagramme d'activité}

Le diagramme d'activité (Figure~\ref{fig:activity_diagram}) modélise le cycle de vie d'un projet de construction, depuis
sa création jusqu'à sa clôture. Le flux principal comprend : création du projet et
affectation du client, constitution de l'équipe et définition des phases, planification
des tâches, établissement du devis, signature du contrat, exécution des travaux avec
suivi d'avancement, facturation et clôture. Les points de décision portent sur
l'acceptation du devis et l'achèvement du projet ; en cas de refus, le devis peut
être révisé avant reprise du flux.

\vspace{0.4em}
\begin{center}
  \includegraphics[height=0.50\textheight,keepaspectratio]{activiteUML.png}
  \captionof{figure}{Diagramme d'activité : cycle de vie d'un projet de construction}
  \label{fig:activity_diagram}
\end{center}

\newpage
\section*{Conclusion}

Ce chapitre a formalisé la conception de l'application de gestion de projets de
CIVCO BTP à travers trois niveaux d'abstraction complémentaires.

\medskip

Au niveau \textbf{fonctionnel}, l'analyse des besoins a identifié six acteurs métier,
vingt cas d'utilisation et treize exigences fonctionnelles couvrant l'ensemble du
cycle de vie d'un projet de construction : de la relation client à la facturation,
en passant par la planification des tâches et le suivi contractuel. Les neuf
exigences non-fonctionnelles fixent les contraintes de performance, de sécurité,
d'ergonomie et de maintenabilité auxquelles la solution doit répondre.

\medskip

Au niveau \textbf{architectural}, l'application repose sur une séparation nette en
trois couches : interface React en SPA, API Laravel exposant des services REST
sécurisés par Sanctum et RBAC, et persistance relationnelle via Eloquent. Le flux
d'authentification, la matrice des permissions et l'organisation modulaire du code
(backend par contrôleurs/services, frontend par features) garantissent une base
solide pour l'évolution future de l'outil.

\medskip

Au niveau \textbf{métier}, la conception détaillée des sept modules (projets, clients,
documents commerciaux, contrats, tâches, tableau de bord, portail client) traduit
les processus internes de CIVCO BTP en composants logiciels cohérents, avec des
règles de gestion explicites et des flux documentaires tracés.

\medskip

Enfin, la modélisation \textbf{UML} --- cas d'utilisation, séquence, classes et
activité --- formalise les interactions, le modèle de données et les flux métier,
et sert de référence pour la phase de réalisation décrite au chapitre suivant.

\medskip

Le chapitre suivant décrit l'implémentation concrète de cette conception : choix
d'environnement, structure du code, captures d'écran des modules principaux et
résultats de la campagne de tests.

\newpage
